% Reviewer-facing sections included by the generated submission source.
\section{Robustness, scope, and method boundaries}
\subsection{Boundary with value of information and goal-oriented design}
Bayesian experimental design selects observations by expected utility, with information gain as one commonly used utility \cite{ChalonerVerdinelli1995}. Targeted and goal-oriented optimal experimental design already select experiments to reduce uncertainty in a specified prediction or quantity of interest rather than the full parameter/state \cite{VanlierEtAl2012,AttiaEtAl2018,ZhongEtAl2026}. Ecological value-of-information analysis values data through their expected improvement in management outcomes \cite{CanessaEtAl2015}. Target-safe design is compatible with all of these decision-theoretic traditions: its reporting rule and experiment cost can be embedded in a constrained utility or loss formulation. Its contribution is therefore not target orientation itself, nor a claim that value of information is incapable of targeting decisions. It is an explicit finite reportability layer linking compatible worlds, target-preserving abstraction, abstention, observation-failure architecture, and a false-resolution contract.

The distinction is operational. Full-world information gain rewards reduction anywhere in the latent state. Goal-oriented OED instead concentrates information on a predictive quantity of interest; management VOI values information through management consequences. Target-safe design may recommend the same experiment as either targeted approach. It additionally asks whether the realized record supports a singleton target within a stated error contract, permits a sharp set-valued report when it does not, and constructs the coarsest action-stable latent quotient preserving that report.

\begin{table}[ht]
\centering
\small
\begin{tabular}{p{0.20\linewidth}p{0.25\linewidth}p{0.45\linewidth}}
\toprule
Approach & Design objective & Relation to target-safe modelling \\
\midrule
Full-world information gain & Reduce entropy over the latent state & Can prioritize precisely measured distinctions that do not alter the declared prediction. \\
Goal-oriented / targeted OED & Reduce uncertainty or maximize information in a prediction/QoI & Already target-relevant; does not by itself define Paper B's compatible-world quotient, sharp abstention report, or failure-aware reporting contract. \\
Bayesian experimental design & Maximize a declared expected utility & General umbrella that can encode target-safe loss and broader objectives. \\
Management value of information & Improve expected management outcome & Directly decision-centred; can recommend the same information when management utilities are specified. \\
Target-safe design & Satisfy a prediction-reporting error contract at low cost & Makes compatible worlds, set-valued abstention, target-preserving state reduction, and observation-failure constraints explicit. \\
\bottomrule
\end{tabular}
\caption{Target-safe modelling is positioned as a finite reportability layer alongside, rather than a replacement for, target-oriented and decision-theoretic experimental design.}
\label{tab:method-boundaries}
\end{table}

\subsection{Dependence on the declared prediction target}
The target-safe quotient is intentionally target-relative. This is not a nuisance parameter of the method but the scientific question supplied by the analyst. Holding the 16 worlds, prior, costs, and likelihood kernels fixed, changing the target from management response to context classification reverses the selected experiment (Table~\ref{tab:target-switch}). Thus no experiment is intrinsically relevant or irrelevant; relevance is defined by whether it separates worlds with different values of the declared target. This target dependence is consistent with targeted/goal-oriented OED; the additional Paper B question is what the resulting finite record can defensibly report.

\begin{table}[ht]
\centering
\small
\input{generated/paper_b_target_switch.tex}
\caption{Probability that each experiment produces a contract-valid deterministic report under two declared targets. The same evidence model yields opposite target-safe choices.}
\label{tab:target-switch}
\end{table}

\subsection{Sensitivity to the false-resolution contract}
The 5\% limit is a baseline reporting contract, not a property hard-coded into the framework. Table~\ref{tab:threshold-sensitivity} varies the limit while holding state sensitivity at 0.95, response accuracy at 0.95, independent screening, and common-failure probability at zero. A stricter 1\% contract rejects the response experiment as insufficiently reliable and therefore retains more ambiguity and spends less. At 5\% and 10\%, the response experiment is admissible. This discontinuity is expected: the method treats the declared error limit as a constraint rather than as a smooth preference weight.

\begin{table}[ht]
\centering
\small
\input{generated/paper_b_threshold_sensitivity.tex}
\caption{Target-safe terminal reports under alternative false-resolution limits. Wrong deterministic reporting remains below the declared limit, while ambiguity and cost respond to the contract.}
\label{tab:threshold-sensitivity}
\end{table}

\subsection{Finite representation and continuous ecological models}
The exact theorems apply to a declared finite world set. This limitation should not be hidden. Continuous-state occupancy, demographic, ecosystem, Bayesian, particle-filter, and ensemble models may supply a finite approximation through posterior draws, particles, scenarios, parameter bins, or representative model members. Target-safe conclusions are then exact for that finite representation and conditional on its support. Refining the representation can reveal target distinctions absent from a coarse approximation, so discretization adequacy must be checked against the prediction target. We do not claim a universal continuous-state convergence theorem.

\subsection{Posterior-sample bridge demonstration}
To make that connection executable, we generated deterministic posterior draws from a continuous invasion-control model with uncertain control effect, spread rate, and source pathway. This is a methodological demonstration rather than an empirical dataset. Each posterior draw is treated as a finite world, and the prediction target is whether eradication or containment is supported. A response assay depends on the continuous control-effect and spread parameters, whereas source tracing identifies a four-level pathway variable that does not alter the management target. Across 500, 2,000, and 10,000 posterior draws, target-safe design consistently selects the response assay; the source-tracing experiment never produces a contract-valid management report (Table~\ref{tab:posterior-bridge}). The selected experiment is therefore stable as the posterior support is refined in this example, while the conclusion remains explicitly conditional on the generated posterior and target definition.

\input{generated/paper_b_posterior_bridge_table.tex}

\subsection{Example of a changed ecological action}
Consider an invasive population that has been detected. If the declared target is current occurrence, further response experiments are unnecessary. If the target is whether eradication is likely to succeed, worlds with equal occurrence but different demographic response must remain distinct, and an intervention-response experiment can be selected. If the target instead concerns source region for pathway regulation, a context or provenance measurement may become decisive. The framework therefore does not merely recommend ``more monitoring''; it changes which observation is justified by the management-facing prediction. This follows the broader principle that conservation monitoring should be embedded in scientific or management objectives rather than treated as stand-alone surveillance \cite{NicholsWilliams2006}.
